\documentclass{article}

\usepackage[a4paper,margin=28mm]{geometry}
\usepackage{osgstyler-lua}

% Resource palettes are independent of the document class and of individual
% formatting roles.
\DeclareOsgFontPalette{editorial}{
  main      = {font = {Latin Modern Roman}},
  companion = {font = {Latin Modern Sans}},
  display   = {alias = companion},
  accent    = {alias = companion},
  mono      = {font = {Latin Modern Mono}},
  symbol    = {alias = main}
}

\DeclareOsgColorPalette{editorial}{
  primary    = {RGB}{36,74,115},
  secondary  = {RGB}{79,109,122},
  accent     = {RGB}{217,119,6},
  text       = {RGB}{32,36,40},
  background = {RGB}{255,253,248},
  structure  = {alias}{primary},
  alert      = {RGB}{166,27,27},
  example    = {RGB}{40,122,80}
}

\DeclareOsgMetricPalette{editorial}{
  spacing = {xs = {.15em}, sm = {.35em}},
  padding = {compact = {.22em}, normal = {.45em}},
  rule    = {thin = {.45pt}, normal = {.75pt}},
  radius  = {small = {2pt}}
}

% Decorations describe visual recipes. Several named instances of the same
% family can be active at once.
\DeclareOsgDecoration{double-line}{
  line = {
    under = {
      position  = below,
      pattern   = solid,
      color     = primary,
      thickness = rule/thin,
      offset    = .12ex
    },
    over = {
      position  = above,
      pattern   = solid,
      color     = primary,
      thickness = rule/thin,
      offset    = .12ex
    }
  }
}

\DeclareOsgDecoration{term-marker}{
  background = {
    color   = accent,
    opacity = 1,
    shape   = text
  }
}

\DeclareOsgDecoration{inline-note}{
  enclosure = {
    border-color = secondary,
    fill-color   = background,
    pattern      = solid,
    thickness    = rule/thin,
    radius       = 0pt,
    padding      = padding/compact
  },
  affix = {
    label = {
      position = before,
      content  = {i},
      color    = primary,
      font     = accent,
      gap      = spacing/sm
    }
  }
}

% Inline template instances combine typography, colors, decorations, and a
% semantic role. The Lua backend renders the decoration; semantic processing
% remains independently replaceable.
\DeclareOsgInlineStyle{term}{
  font       = none,
  weight     = bold,
  shape      = inherit,
  color      = text,
  decoration = term-marker,
  semantics  = none
}

\DeclareOsgInlineStyle{doubly-emphasized}{
  font       = none,
  weight     = inherit,
  shape      = inherit,
  color      = primary,
  decoration = double-line,
  semantics  = emphasis
}

\DeclareOsgInlineStyle{note}{
  font       = none,
  weight     = inherit,
  shape      = inherit,
  color      = text,
  decoration = inline-note,
  semantics  = none
}

% The predefined instances can be themed without changing command bindings.
\EditOsgInlineStyle{emphasis}{color = accent, shape = italic}
\EditOsgInlineStyle{strong}{color = alert, weight = bold}

% Existing and newly introduced author commands all use one-input sockets.
\UseOsgStandardInlineCommands
\BindOsgInlineCommand{\term}{term}
\BindOsgInlineCommand{\doubleemph}{doubly-emphasized}
\BindOsgInlineCommand{\inlinenote}{note}

\UseOsgFontPalette{editorial}
\UseOsgColorPalette{editorial}
\UseOsgMetricPalette{editorial}

\title{An \texttt{article} styled with \texttt{osgstyler}}
\author{osglecture example}
\date{}

\begin{document}
\maketitle

\section{Standard author commands}

The familiar commands now select template instances: \emph{emphasis keeps its
semantic role but receives the palette accent}, while \textbf{strong text uses
the alert color}. The visual commands \uline{underline text} and
\highlight{highlighted text} are rendered by the Lua backend.

The bindings are ordinary sockets. They can be returned locally to their
captured LaTeX defaults:

\begin{quote}
  \UseOsgDefaultInlineCommands
  Here \emph{emphasis} and \textbf{bold text} use the original definitions;
  commands introduced by \texttt{osgstyler}, such as
  \uline{underline} and \highlight{highlight}, become transparent.
\end{quote}

Outside the quotation, the surrounding socket assignments are active again:
\emph{accented emphasis}, \textbf{alert strong text}, and
\highlight{a visible marker}.

\section{Custom instances and commands}

A project-specific command can be bound without defining its implementation
twice. A \term{key term} combines bold text and a palette-based marker.
The command can \doubleemph{combine overlining and underlining} in one style.

The standard command \fbox{uses the palette's primary colour}. An enclosure
and an affix can also be composed:
\inlinenote{a short, unbreakable framed note}.

\section{Resource switching}

The styles above contain no literal font names and almost no literal
dimensions. Fonts come from the font palette, colors from the color palette,
and reusable dimensions from the metric palette. A theme can therefore replace
the resources without redefining \verb|\term|, \verb|\doubleemph|, or any
template instance.

\end{document}
